\chapter{Last-FM* Experimental Results and Ablation Analysis}
\label{ch:lastfm-benchmark}

This chapter reports controlled Last-FM* experiments organized by scientific question:
structural capacity, contribution of explicit candidate context, routing criterion, and full-catalog seed stability.
All numbers are development validation results on frozen checkpoints.
The sealed external test is outside the scope of this chapter.

\section{HGT Backbone Selection}
\label{sec:lastfm-backbone}

To select the structural capacity under the common development protocol, HGT width, depth and number of attention heads were compared while holding the decoder family and training protocol fixed.
To isolate backbone capacity from explicit-context effects, selection used the shared sampled-validation checkpoint metric before freezing the explicit-context ablations.

The selected configuration is
\begin{equation}
d=64,\qquad L=2,\qquad H=2,
\end{equation}
with dropout $0.1$.
Larger width, deeper stacks or additional heads did not justify replacing this setting under the frozen development protocol.

\section{Contribution of Explicit Candidate Context}
\label{sec:lastfm-ablation}

With the HGT backbone fixed, the controlled ablation ladder isolates the contribution of explicit candidate context:
\begin{enumerate}
\item HGT only (257-dimensional structural input);
\item HGT + A5 (262-dimensional main input);
\item HGT + A5 + H3 (265-dimensional main input);
\item TRUE FINAL (265-dimensional main input plus the LEG$_{K2}$ residual).
\end{enumerate}
All variants use the same HGT capacity, the same frozen training candidate tables, the same optimizer family and the seeds $\{101,202,303\}$.

\begin{table}[htbp]
\centering
\footnotesize
\caption{Progressive addition of explicit candidate context.
Sampled NDCG@20 is the checkpoint-selection metric; full-catalog NDCG@20 and Recall@20 evaluate the same frozen checkpoints on the eligible catalog.}
\label{tab:lastfm-ablation}
\setlength{\tabcolsep}{2.6pt}
\begin{tabular}{@{}lcccc@{}}
\toprule
Variant & Main decoder & Sampled NDCG@20 & Full-cat.\ NDCG@20 & Full-cat.\ Recall@20 \\
\midrule
HGT only & 257D & \(0.4755\) & \(0.0093\) & \(0.0169\) \\
HGT + A5 & 262D & \(0.7908\) & \(0.0640\) & \(0.1128\) \\
HGT + A5 + H3 & 265D & \(0.8729\) & \(0.2321\) & \(0.3647\) \\
TRUE FINAL (+ LEG) & 265D + residual & \(0.8733\) & \(0.2764\) & \(0.3768\) \\
\bottomrule
\end{tabular}
\end{table}

Per-seed sampled NDCG@20 values were:
HGT only \(0.475024\), \(0.475407\), \(0.475953\);
HGT+A5 \(0.791067\), \(0.790560\), \(0.790768\);
HGT+A5+H3 \(0.872794\), \(0.872930\), \(0.873007\);
TRUE FINAL \(0.873224\), \(0.873253\), \(0.873515\).

A5 introduces the first large improvement on both protocols.
H3 produces the largest full-catalog increase.
LEG changes sampled NDCG@20 by only about $0.0004$ relative to HGT+A5+H3, while adding a further full-catalog gain of about $0.044$ NDCG@20.

\begin{figure}[htbp]
\centering
\includegraphics[width=\textwidth]{fig_lastfm_ablation_bars.pdf}
\caption{Last-FM* ablation ladder.
Panel~(a): sampled validation NDCG@20 used for checkpoint selection.
Panel~(b): full-catalog NDCG@20 of the same frozen checkpoints.
Error bars show the seed standard deviation over three training seeds.
The panels use different $y$-axis scales because they evaluate different candidate universes under the same NDCG@20 formula.}
\label{fig:lastfm-ablation-bars}
\end{figure}

\section{Routing Criterion Ablation}
\label{sec:lastfm-routing}

The routing stage chooses which history items enter the Top25 local context.
After a Top25 set is chosen, the selected items are represented using signed A11/Pearson $\phi$ values in H3 and LEG$_{K2}$.
Alternative routing criteria therefore change only the selection rule, not the downstream H3 construction.

To identify the most useful candidate-conditioned history context, the following item--pair association criteria were compared under a fixed Top25 budget:
\begin{align}
J_{\mathrm{pair}}(h,X)
&=
\frac{n_{11}}{n_h+n_X-n_{11}},
\label{eq:lastfm-Jpair}\\
C_{\mathrm{pair}}(h,X)
&=
\frac{n_{11}}{\sqrt{n_h n_X}},
\label{eq:lastfm-Cpair}
\end{align}
together with raw co-occurrence $n_{11}$, mutual information (MI), normalized pointwise mutual information (NPMI), and the unsigned $G^{2}$ likelihood-ratio statistic.
These quantities are distinct from the set-level A5 overlaps $J_{\mathrm{set}}$ and $C_{\mathrm{set}}$ of Section~\ref{sec:lastfm-A5}.

Held fixed across routing variants:
the candidate pairs, the Top25 budget $K=\min(25,|H_u\setminus\{X\}|)$, the downstream signed-A11 H3/LEG construction, and the main architecture.
Only the routing criterion changes.

The frozen routing policy retained for TRUE FINAL is signed A11/Pearson $\phi$.
Quantitative full-catalog tables for every routing alternative are not included here because corresponding frozen result artefacts were not located in the sealed development archive; the methodological comparison and the frozen winner are stated above.

\section{Full-Catalog Results and Seed Stability}
\label{sec:lastfm-fullrank}

Epochs were selected using sampled validation NDCG@20.
Full-catalog evaluation was then applied to the resulting frozen checkpoints without re-selection.
All full-ranking runs used 23{,}163 validation users with at least one positive.

\begin{table}[htbp]
\centering
\footnotesize
\caption{TRUE FINAL full-catalog metrics by seed.
Sampled NDCG@20 selected the epoch; full-catalog metrics are evaluation only.}
\label{tab:lastfm-true-final}
\setlength{\tabcolsep}{3pt}
\begin{tabular}{@{}crrrrr@{}}
\toprule
Seed & Best epoch & Sampled NDCG@20 & Full-cat.\ NDCG@20 & Recall@20 & MRR \\
\midrule
101 & 36 & 0.873224 & 0.277432 & 0.375599 & 0.356988 \\
202 & 34 & 0.873253 & 0.279879 & 0.378722 & 0.359681 \\
303 & 43 & 0.873515 & 0.271886 & 0.376197 & 0.343391 \\
\midrule
Mean \(\pm\) SD & --- & \(0.8733\pm 0.0002\) & \(0.2764\pm 0.0041\) & \(0.3768\) & \(0.3534\) \\
\bottomrule
\end{tabular}
\end{table}

\begin{table}[htbp]
\centering
\small
\caption{Full-catalog NDCG@20 by seed for the ablation layers without LEG.}
\label{tab:lastfm-fullrank-ablation}
\setlength{\tabcolsep}{4pt}
\begin{tabular}{crrr}
\toprule
Seed & HGT only & HGT+A5 & HGT+A5+H3 \\
\midrule
101 & 0.008548 & 0.064571 & 0.227969 \\
202 & 0.009520 & 0.060866 & 0.241332 \\
303 & 0.009935 & 0.066630 & 0.227054 \\
\midrule
Mean & \(0.0093\) & \(0.0640\) & \(0.2321\) \\
\bottomrule
\end{tabular}
\end{table}

For TRUE FINAL, mean full-catalog NDCG values were additionally $\mathrm{NDCG}@5=0.2336$ and $\mathrm{NDCG}@10=0.2503$; HitRate@20 was approximately $0.664$.

Under the full-catalog development protocol, the controlled ablations support three factual conclusions:
A5 supplies the first large ranking gain;
H3 supplies the largest full-catalog gain among the explicit-context components;
LEG$_{K2}$ leaves the sampled checkpoint metric nearly unchanged while providing an additional full-catalog refinement.
External baseline head-to-head numbers are omitted until the sealed external evaluation protocol is completed and approved for inclusion.
